% from ega0-1-fr.tex
\oldpage[0\textsubscript{I}]{39}
le long des $V_{\lambda\mu}$~; si on identifie $X_\lambda$ à la partie ouverte
$X'_\lambda$ correspondante dans $X$, les hypothèses entraînent que les trois
ensembles $V_{\lambda\mu}\cap V_{\lambda\nu}$,
$V_{\mu\nu}\cap V_{\mu\lambda}$,
$V_{\nu\lambda}\cap V_{\nu\mu}$ s'identifient à
$X'_\lambda\cap X'_\mu\cap X'_\nu$. On peut alors transporter à
$X'_\lambda$ la structure d'espace annelé de $X_\lambda$, et si
$\mathcal{A}'_\lambda$ est le faisceau d'anneaux transporté de
$\mathcal{A}_\lambda$, les $\mathcal{A}'_\lambda$ vérifient la condition de
recollement (3.3.1) et définissent donc un faisceau d'anneaux $\mathcal{A}$ sur
$X$~; on dit que $(X,\mathcal{A})$ est l'espace annelé obtenu par
\emph{recollement des $(X_\lambda,\mathcal{A}_\lambda)$ le long des
$V_{\lambda\mu}$}, au moyen des $\varphi_{\lambda\mu}$.
\end{env}

\subsection{Image directe d'un $\mathcal{A}$-Module}
\label{subsection:0.4.2-fr}

\begin{env}[4.2.1]
\label{0.4.2.1-fr}
Soient $(X,\mathcal{A})$, $(Y,\mathcal{B})$ deux espaces annelés,
$\Psi=(\psi,\theta)$ un morphisme
$(X,\mathcal{A})\to(Y,\mathcal{B})$~; $\psi_*(\mathcal{A})$ est donc un
faisceau d'anneaux sur $Y$, et $\theta$ un homomorphisme
$\mathcal{B}\to\psi_*(\mathcal{A})$ de faisceaux d'anneaux. Soit alors
$\mathcal{F}$ un $\mathcal{A}$-Module~; son image directe
$\psi_*(\mathcal{F})$ est un faisceau de groupes abéliens sur $Y$. En outre,
pour tout ouvert $U\subset Y$,
\[
  \Gamma(U,\psi_*(\mathcal{F}))
  =\Gamma(\psi^{-1}(U),\mathcal{F})
\]
est muni d'une structure de module par rapport à l'anneau
$\Gamma(U,\psi_*(\mathcal{A}))=\Gamma(\psi^{-1}(U),\mathcal{A})$~; les
applications bilinéaires qui définissent ces structures étant compatibles
avec les opérations de restriction, définissent sur $\psi_*(\mathcal{F})$ une
structure de $\psi_*(\mathcal{A})$-Module. L'homomorphisme
$\theta:\mathcal{B}\to\psi_*(\mathcal{A})$ permet alors de définir aussi sur
$\psi_*(\mathcal{F})$ une structure de $\mathcal{B}$-Module~; nous dirons que
ce $\mathcal{B}$-Module est l'\emph{image directe de $\mathcal{F}$ par le
morphisme $\Psi$}, et nous le noterons $\Psi_*(\mathcal{F})$. Si
$\mathcal{F}_1$, $\mathcal{F}_2$ sont deux $\mathcal{A}$-Modules sur $X$ et
$u$ un $\mathcal{A}$-homomorphisme $\mathcal{F}_1\to\mathcal{F}_2$, il est
immédiat (en considérant les sections au-dessus des ouverts de $Y$) que
$\psi_*(u)$ est un $\psi_*(\mathcal{A})$-homomorphisme
$\psi_*(\mathcal{F}_1)\to\psi_*(\mathcal{F}_2)$, et \emph{a fortiori} un
$\mathcal{B}$-homomorphisme
$\Psi_*(\mathcal{F}_1)\to\Psi_*(\mathcal{F}_2)$~; en tant que
$\mathcal{B}$-homomorphisme, nous le noterons $\Psi_*(u)$. On voit donc que
$\Psi_*$ est un \emph{foncteur covariant} de la catégorie des
$\mathcal{A}$-Modules dans celle des $\mathcal{B}$-Modules. En outre, il est
immédiat que ce foncteur est \emph{exact à gauche} (G, II, 2.12).

Sur $\psi_*(\mathcal{A})$, la structure de $\mathcal{B}$-Module et la
structure de faisceaux d'anneaux définissent une structure de
$\mathcal{B}$-Algèbre~; on notera $\Psi_*(\mathcal{A})$ cette
$\mathcal{B}$-Algèbre.
\end{env}

\begin{env}[4.2.2]
\label{0.4.2.2-fr}
Soient $\mathcal{M}$, $\mathcal{N}$ deux $\mathcal{A}$-Modules. Pour tout
ouvert $U$ de $Y$, on a une application canonique
\[
  \Gamma(\psi^{-1}(U),\mathcal{M})
  \times\Gamma(\psi^{-1}(U),\mathcal{N})
  \longrightarrow
  \Gamma(\psi^{-1}(U),\mathcal{M}\otimes_{\mathcal{A}}\mathcal{N})
\]
qui est bilinéaire sur l'anneau
$\Gamma(\psi^{-1}(U),\mathcal{A})=\Gamma(U,\psi_*(\mathcal{A}))$, et
\emph{a fortiori} sur $\Gamma(U,\mathcal{B})$~; elle définit donc un
homomorphisme
\[
  \Gamma(U,\Psi_*(\mathcal{M}))
  \otimes_{\Gamma(U,\mathcal{B})}
  \Gamma(U,\Psi_*(\mathcal{N}))
  \longrightarrow
  \Gamma(U,\Psi_*(\mathcal{M}\otimes_{\mathcal{A}}\mathcal{N}))
\]
et comme on vérifie aussitôt que ces homomorphismes sont compatibles avec les
opérations de restriction, ils donnent un homomorphisme canonique fonctoriel
de $\mathcal{B}$-Modules
\[
  \Psi_*(\mathcal{M})\otimes_{\mathcal{B}}\Psi_*(\mathcal{N})
  \longrightarrow
  \Psi_*(\mathcal{M}\otimes_{\mathcal{A}}\mathcal{N}).
  \tag{4.2.2.1}\label{0.4.2.2.1-fr}
\]

